\section{方法对一般复合算符的应用}
\label{sec:casivari}

上一节以标量密度与赝标密度的重正化为原型给出了非微扰方案的一般策略。本节简要描述其他情形下所需
的必要修改。我们将遵循第 \ref{sec:idea} 节引入的算符分类。

\subsection{有限算符}

对矢量流与轴矢流，我们没有任意施加重正化条件的自由：重正化条件只有在与 Ward 恒等式相容时才是
可接受的。具体地，可以对局域矢量流应用式 (\ref{eq:zdef}) 来确定 $Z_{V^L}$：
\beq Z_{V^L} = \frac{Z_\psi}{\Gamma_{V^L}}, \label{eq:zvl} \eeq
我们现在证明，这个由上一节一般处方得到的 $Z_V^L$ 的定义，与从夸克态上矢量流 Ward 恒等式导出的定义
相一致 \cite{bo,ka}\footnote{在本小节余下部分我们在格点单位下工作，取 $a=1$。}：
\beq
q_{\mu} \Bigl[Z_{V^L} \Lambda_{V^L_\mu}(p+\frac{q}{2},p-\frac{q}{2})
\Bigr]=-i\Bigl( S^{-1}(p+\frac{q}{2}) - S^{-1}(p-\frac{q}{2})\Bigr).
\label{eq:pallino}\eeq
其中 $\Lambda_{V^L_\mu}(p+q/2,p-q/2)$ 是带动量转移 $q$ 的截腿顶点。在式 (\ref{eq:pallino}) 两边施加
$q=0$ 处的 $\partial / \partial q_\rho$，得到
\beq
Z_{V^L}\Bigl( \Lambda_{V^L_\rho}(p) + q_\mu \frac{\partial}{\partial q_\rho}
 \Lambda_{V^L_\mu}(p+\frac{q}{2},p-\frac{q}{2})\vert_{q=0}\Bigr)
=-i \frac{\partial}{\partial p_\rho}S^{-1}(p).
\label{eq:pinco} \eeq
左端第二项在 $q \to 0$ 时消失。用 $\gamma_{\rho}$ 对式 (\ref{eq:pinco}) 求迹，
便恢复了式 (\ref{eq:zvl}) 中 $Z_{V^L}$ 的定义。

人们会想对局域轴矢流走同样的路径，即定义
\beq Z_{A^L} = \frac{Z_\psi}{\Gamma_{A^L}}.  \label{eq:zal} \eeq
然而 $Z_{A^L}$ 的这一定义并不等价于由 Ward 恒等式
\beq
q_{\mu} \Bigl[Z_{A^L} \Lambda_{A^L_\mu}(p+\frac{q}{2},p-\frac{q}{2})
\Bigr]=-i\Bigl( \gamma_5 S^{-1}(p+\frac{q}{2}) + S^{-1}(p-\frac{q}{2})
 \gamma_5 \Bigr).
\label{eq:awi}\eeq
所得的定义。为看清这一点，我们对式 (\ref{eq:awi}) 两边施加 $\partial /\partial
 q_\rho$，如矢量流的情形那样：
\beq Z_{A^L}\Bigl( \Lambda_{A^L_\rho}(p) + q_\mu
\frac{\partial}{\partial q_\rho}
 \Lambda_{A^L_\mu}(p+\frac{q}{2},p-\frac{q}{2})\vert_{q=0}\Bigr) =
  \nn \eeq
\beq  -\frac{i}{2}\Bigl( \gamma_5 \frac{\partial}{\partial p_\rho}S^{-1}(p)
-\frac{\partial}{\partial p_\rho}S^{-1}(p) \gamma_5 \Bigr).
\label{eq:dawi} \eeq
但在轴矢情形，左端第二项当 $q \to 0$ 时并不消失；事实上，在存在无质量 Goldstone 玻色子时，正是这一项
使 Ward 恒等式得以饱和。然而正如附录所证明的，该项在高 $p^2$ 值处变得可以忽略。此时对式
(\ref{eq:dawi}) 求迹便恢复 (\ref{eq:zal})。总之我们预期：在大 $\mu^2$ 处使用式 (\ref{eq:zal})，将得到
接近于由强子态或夸克态上 Ward 恒等式所得的 $Z_{A^L}$ 值 \cite{wi}。

$Z_S/Z_P$ 是另一个有限的量，它完全由 Ward 恒等式确定，尽管 $Z_S$ 与 $Z_P$ 各自的重正化条件是任意
的。换言之，人们可以自由决定两个算符之一（比如说标量密度）的重正化条件，而 $Z_P$ 的值将自动随之
确定。第 \ref{sec:idea} 节的流程尊重 Ward 恒等式，从而保证得到正确的 $Z_S/Z_P$ 值。当然，出于与轴矢流
相同的原因，这只在大 $\mu^2$ 值处成立。

\subsection{对数发散算符}

对这一类算符，只剩下算符混合问题需要讨论。混合有两种：一种是连续统中也可能发生的，另一种是由格点
手征对称性显式破缺诱导的。这里我们只关心后者，因而把讨论限制于在连续统中乘性重正化的算符。为说明
我们的论证，考虑四费米子算符
\beq O^{\Delta S=2} = \Bigl(\bar s \gamma^\mu (1-\gamma_5) d \Bigr)
\Bigl(\bar s \gamma_\mu (1-\gamma_5) d \Bigr), \label{eq:ds2b} \eeq
它与 $K^0$--$\bar K^0$ 混合振幅的计算相关。由于 Wilson 项，该算符与不同手征性的算符混合
\cite{bo,te,4ferm}：
\beq  O^{\Delta S=2}_{{\rm cont}} = Z_O \Bigl( O^{\Delta S=2}_{{\rm
latt}} + \sum_i Z_i O^i \Bigr),\label{eq:omix} \eeq
其中算符 $O^i$ 为 $O^1= (\bar s d )(\bar s d),
O^2=(\bar s
\gamma_5 d)(\bar s \gamma_5 d)$ 等 \cite{4ferm}。混合系数 $Z_i$ 是 $g_0^2(a)$ 的有限函数。整体重正化
$Z_O$ 用于消除对数发散。

现在概述 $O^{\Delta S=2}$ 的重正化流程。先写出 $O^{\Delta S=2}_{{\rm latt}}$ 在 $p^2=\mu^2$
的夸克态之间的四点截腿格林函数
\beq \Lambda_{\alpha \beta \gamma \delta}^{ABCD}(pa),
\label{eq:g4} \eeq
其中 $\alpha, \beta \dots$ 与 $A, B \dots$ 分别是自旋与色指标。为求出混合系数 $Z_i$，如同两夸克算符
那样使用适当的投影算符 $\hat P$\footnote{为简化表述，此处不讨论不同可能色结构带来的复杂性。}，
例如
\beq {\rm tr } ( \Lambda \hat P )= \Lambda_{ \alpha \beta
 \gamma \delta }^{AABB}( \gamma_5)_{ \beta \alpha}
 ( \gamma_5 )_{ \delta \gamma}. \eeq
向对应于算符 $O^i$ 的所有可能色—自旋结构投影，便得到关于混合系数 $Z_i$ 的线性方程组。用这样确定的
$Z_i$，可以计算减除后的算符 $O^s=O^{\Delta S=2}_{{\rm latt}}+ \sum_i Z_i O^i$ 的格林函数。最后，从 $O^s$
的截腿顶点出发，通过施加重正化条件
\beq Z_O \Bigl( \mu , g(a) \Bigr) Z^{-2}_{ \psi }
 \Bigl( \mu , g(a) \Bigr) \Gamma^s_O ( pa) \vert_{p^2 =
 \mu^2} = 1, \eeq
消除整体对数发散，其中
\beq \Gamma^s_O \sim { \rm tr } (
 \Lambda^s \hat P_{ L \times L} ) ,\eeq
而 $P_{ L \times L}$ 是狄拉克张量 $ \gamma_ \mu ( 1 - \gamma_5 ) \otimes \gamma^\mu (
1 - \gamma_5 ) $ 上适当的投影算符。

\subsection{幂次发散算符}

当复合算符可与低维算符混合时，它会以 $1/a$ 发散。例子包括与深度非弹性散射相关的算符，以及出现在
$\Delta S=1/2$ 哈密顿量中的企鹅算符。考虑 $\Delta S=1/2$ 算符
\beq O^{ \Delta S=1/2}= ( \bar s \gamma_\mu ( 1 - \gamma_5 )
d ) \sum_i (\bar q_i \gamma^\mu ( 1 + \gamma_5 ) q_i) .\label{eq:dt12} \eeq
$ O^{ \Delta S=1/2}_{{\rm latt}}$ 通过图 \ref{fig:penguin} 中的图与标量和赝标密度混合（甚至在
$\alpha_s$ 零阶）\footnote{$O^{\Delta S=1/2}_{{\rm latt}}$ 还与色磁算符（例如 $\bar s \sigma_{\mu \nu}
G^{\mu \nu} d$）混合，此处不予考虑。}：
\beq O^{\Delta S=1/2}_{{\rm cont}} =
Z_O O^{\Delta S=1/2}_{{\rm latt}}+ Z_1 (\bar s d )+ Z_2 (\bar s  \gamma_5
 d )+ \dots. \label{eq:ds12} \eeq
\begin{figure}[htbp]
\centering
\includegraphics[width=0.62\textwidth]{images/feynman1.pdf}
\caption{把式 (\ref{eq:dt12}) 中的四费米子算符与标量密度和赝标密度相混合的图。该混合由 Wilson 作用量与
Clover 作用量中存在的手征破缺项诱导。}
\label{fig:penguin}
\end{figure}

为求系数 $Z_{1,2}$，可要求式 (\ref{eq:ds12}) 右端算符的两点格林函数（在离壳夸克态上）为零
\cite{ber1}。一旦去除了所有发散部分，就可以如上讨论的那样施加重正化条件，以去除整体的对数发散。
